\expandafter\def\csname HistoricalNoteText4\endcsname{In a letter to \textsc{Leibnitz}. Horrebowii Opera T. II p. 142.}
\expandafter\def\csname HistoricalNoteText5\endcsname{It is therefore equal to the sine of the Sun’s mean horizontal parallax if the mean
distance of the Earth from the Sun is taken as the unit.}
\clearpage\SourceLeaf{99}%
\pdfbookmark[1]{Remarks on simplifying the calculation of geocentric positions of the planets}{bemerkungen1804}
\begin{center}
{\small SOME REMARKS ON SIMPLIFYING THE CALCULATION
OF THE GEOCENTRIC POSITIONS OF THE PLANETS.\par}
\medskip\hrule\medskip
{\small Monatliche Correspondenz zur Beförderung der Erd- und Himmels-Kunde,\par
edited by Baron \textsc{von Zach}. May 1804.\par}
\medskip\hrule
\end{center}\medskip
Since the invention of pendulum clocks, all our observations of fixed stars,
planets, and comets have been referred directly to the equator rather than to the
ecliptic. Our latest and best star catalogues and charts likewise take right
ascension and declination as their basis instead of longitude and latitude. There is
therefore very often occasion to calculate the geocentric positions of planets and
comets relative to the equator from their heliocentric positions in their orbits.
There would be occasion to do so still more often if astronomers decided to use right
ascensions and declinations throughout the astronomical ephemerides, in place of
the planets’ relatively unhelpful longitudes and latitudes. Right ascensions and
declinations are far more useful in every practical respect. The excellent
\textsc{Römer} already recommended this a hundred years ago\HistoricalNote{4}, and it
becomes quite indispensable, in particular, for the two most recently discovered
planets, which are so difficult to observe and can be picked out from the faint
fixed stars around them only with the help of very detailed star charts. Wider
adoption of another proposal would give occasion for this calculation just as
often. When comparing the observed position
\SourceLeaf{100}%
of a planet or comet with its calculated position, take the observed right
ascension and declination directly as the basis instead of first deriving from
them, as is usually done, a so-called observed longitude and latitude. A
competent judge has already mentioned the advantages of this procedure in
vol. V of the Mon. Corr., p. 594.

From this point of view, the planet’s geocentric longitude and latitude are
merely intermediate quantities for finding its position relative to the equator.
It may therefore count in favour of the proposals above that one can dispense
with this intermediate calculation and even with reducing the heliocentric
position in the orbit to the heliocentric position relative to the ecliptic.
Very simple and flexible formulas, to be developed in this paper, give the
geocentric right ascension and declination directly from the position in the
orbit. Another advantage is the great ease with which this method can take
\textit{parallax} into account even when the planet is outside the meridian
of the place of observation. This is less often necessary, but when it is,
other methods make it far more troublesome.

Pass three mutually perpendicular planes through the centre of the Sun,
the first parallel to the Earth’s equator, the second through the equinoctial
points, and hence the third through the solstitial points. Let $Z$, $Y$, $X$
be the perpendicular distances of the centre of the Earth from these planes,
respectively, and $z$, $y$, $x$ the distances of a planet from those same planes.
Take the distances as positive on the side of the first plane toward the North
Pole, of the second toward the summer solstice, and of the third toward the
vernal equinox. Accordingly, $z-Z$, $y-Y$, $x-X$ will be the perpendicular
distances of the planet, taken in the same way, from three planes parallel to
those above and passing through the centre of the Earth. If, then, $\alpha$
denotes the geocentric right ascension of the planet, $\delta$ its declination,
and $\Delta$ its distance from the Earth, we have
\[
\begin{aligned}
x-X&=\Delta\cos\delta\cos\alpha\\
y-Y&=\Delta\cos\delta\sin\alpha\\
z-Z&=\Delta\sin\delta
\end{aligned}
\]
\SourceLeaf{101}%
Thus $\alpha$ is found from the formula $\operatorname{tang}\alpha=\frac{y-Y}{x-X}$.
The sign of the numerator, positive or negative, must determine whether
$\alpha$ is taken in the first two quadrants or the last two. Then we obtain
\[
\Delta\cos\delta=\frac{x-X}{\cos\alpha}=\frac{y-Y}{\sin\alpha},\quad
\text{and}\quad\operatorname{tang}\delta=\frac{z-Z}{\Delta\cos\delta}
\]

In this way we obtain the planet’s right ascension and declination as seen from the
centre of the Earth. To obtain them as they appear from a point on the Earth’s surface,
the only change needed in the formulas above is to use the distances of the place
of observation from the three fundamental planes in place of the coordinates
$X$, $Y$, $Z$ of the centre. Let the radius of the Earth be $=\rho$\HistoricalNote{5},
the elevation of the pole at the place of observation $=\varphi$, and the sidereal
time there at the instant of observation, expressed as an arc, or the right ascension
of the culminating point of the equator, $=\vartheta$. Then, as is easy to see,
those distances are
\[
\begin{gathered}
X+\rho\cos\varphi\cos\vartheta\\
Y+\rho\cos\varphi\sin\vartheta\\
Z+\rho\sin\varphi
\end{gathered}
\]
Here the Earth is treated as a sphere. If it were found necessary to take the
spheroidal shape of the Earth into account as well, as might happen with comets
that come very close to the Earth, one need only substitute for $\rho$ the distance
from the place of observation to the Earth’s centre, and for $\varphi$ its
so-called corrected elevation of the pole. These are determined by known rules.

We now see that all that remains is to find a convenient method of determining
the coordinates $X$, $Y$, $Z$, $x$, $y$, $z$. For this purpose, describe a sphere
of unspecified radius around the Sun. On this sphere let $P$ denote the north pole
of the ecliptic, $p$ the north pole of the plane of the planet’s orbit, $K$ the
position of the Earth, $k$ the heliocentric position of the planet, and finally
$\mathfrak{x}$, $\mathfrak{y}$, $\mathfrak{z}$ the poles of the three fundamental planes
on the sides where the distances $x$, $y$, $z$ are taken as positive.
\SourceLeaf{102}%
Thus $\mathfrak{z}$ denotes the north pole of the equator, $\mathfrak{x}$ the point
of the vernal equinox, and $\mathfrak{y}$ the point on the equator whose right
ascension is $90^0$. Anyone who needs a figure will find it easy to draw one from
this description. Now, if we put the distance of the Earth from the Sun $=R$,
we clearly have
\[
\begin{aligned}
X&=R\cos\mathfrak{x}K\\
Y&=R\cos\mathfrak{y}K\\
Z&=R\cos\mathfrak{z}K
\end{aligned}
\]
Thus, since in the spherical triangle $\mathfrak{x}PK$ the side $PK=90^0$,
and hence $\cos\mathfrak{x}K=\sin\mathfrak{x}P\cos\mathfrak{x}PK$, we have
\[
\begin{aligned}
X&=R\sin\mathfrak{x}P\cos\mathfrak{x}PK,\quad\text{and likewise}\\
Y&=R\sin\mathfrak{y}P\cos\mathfrak{y}PK\quad\text{and}\\
Z&=R\sin\mathfrak{z}P\cos\mathfrak{z}PK
\end{aligned}
\]

In just the same way, if we denote the planet’s distance from the Sun by $r$,
its coordinates will be
\[
\begin{aligned}
x&=r\sin\mathfrak{x}p\cos\mathfrak{x}pk\\
y&=r\sin\mathfrak{y}p\cos\mathfrak{y}pk\\
z&=r\sin\mathfrak{z}p\cos\mathfrak{z}pk
\end{aligned}
\]

We stipulate here once and for all that the spherical angle $\mathfrak{x}PK$
is to be understood in the sense in which the ray $PK$ comes after the ray
$P\mathfrak{x}$, \textit{following the order of the symbols}. Thus it is distinct from
$KP\mathfrak{x}$, and together the two angles make $360^0$. Every other
spherical angle is to be understood in the same way. This more precise
convention has the advantage that the basic formulas of spherical trigonometry
extend directly to triangles with angles greater than $180^0$. We thereby avoid
having to distinguish several separate cases. Throughout, angles whose
difference is $360^0$ or a multiple of it are regarded as equivalent.

We now turn first to the coordinates $X$, $Y$, $Z$. Put the obliquity of the
ecliptic $=\varepsilon$ and the heliocentric longitude of the Earth $=\lambda$
($=$ geocentric longitude of the Sun $+180^0$). In the formulas above we thus have
$\mathfrak{x}P=90^0$, $\mathfrak{y}P=90^0+\varepsilon$, $\mathfrak{z}P=\varepsilon$,
$\mathfrak{x}PK=\lambda$, $\mathfrak{y}PK=\mathfrak{z}PK=\lambda-90^0$, and hence
\SourceLeaf{103}%
\[
\begin{aligned}
X&=R\cos\lambda\\
Y&=R\sin\lambda\cos\varepsilon\\
Z&=R\sin\lambda\sin\varepsilon
\end{aligned}
\]

For the planet we put, for brevity, $\mathfrak{x}p=a$, $\mathfrak{y}p=b$,
$\mathfrak{z}p=c$, its distance along its orbit from the ascending node on the
ecliptic $=t$, and the angles $\mathfrak{x}pP$, $\mathfrak{y}pP$, $\mathfrak{z}pP$
respectively $=A$, $B$, $C$. It is easy to see that $Ppk=t-90^0$ (or, by the
remark above, $=t+270^0$), and so $\mathfrak{x}pk=A+t-90^0$,
$\mathfrak{y}pk=B+t-90^0$, $\mathfrak{z}pk=C+t-90^0$. We thus obtain
\[
\begin{aligned}
x&=r\sin a\sin(A+t)\\
y&=r\sin b\sin(B+t)\\
z&=r\sin c\sin(C+t)
\end{aligned}
\]

It remains to derive the quantities $a$, $A$, and so on, which depend only
on the orientation of the planet’s orbit and are independent of its position
in that orbit at any particular time. They are to be derived from the inclination
of the orbital plane and the longitude of the ascending node. We denote the
former by $i$, the latter by $\AscendingNode$. The triangle $\mathfrak{x}pP$
gives us the following three equations
\[
\begin{aligned}
\operatorname{cotang}\mathfrak{x}pP&=\frac{\sin pP\operatorname{cotang}\mathfrak{x}P-\cos pP\cos pP\mathfrak{x}}{\sin pP\mathfrak{x}}\\
\cos\mathfrak{x}p&=\cos pP\cos\mathfrak{x}P+\sin pP\sin\mathfrak{x}P\cos pP\mathfrak{x}\\
\sin\mathfrak{x}p&=\frac{\sin\mathfrak{x}P\sin pP\mathfrak{x}}{\sin\mathfrak{x}pP}
\end{aligned}
\]

Likewise, the triangles $\mathfrak{y}pP$, $\mathfrak{z}pP$ each give three
similar equations. There is no need to write them out here, since they are
obtained from the three above simply by replacing $\mathfrak{x}$ with
$\mathfrak{y}$ and $\mathfrak{z}$, respectively. Now we have
\[
pP=i,\quad pP\mathfrak{x}=90^0-\AscendingNode,\quad pP\mathfrak{y}=pP\mathfrak{z}=180^0-\AscendingNode
\]
With these and the other substitutions, our nine equations become
\[
\begin{aligned}
\operatorname{cotang}A&=-\cos i\operatorname{tang}\AscendingNode\\
\cos a&=\sin i\sin\AscendingNode\\
\sin a&=\frac{\cos\AscendingNode}{\sin A}
\end{aligned}
\]
\SourceLeaf{104}%
\[
\begin{aligned}
\operatorname{cotang}B&=\frac{-\sin i\operatorname{tang}\varepsilon+\cos i\cos\AscendingNode}{\sin\AscendingNode}\\
\cos b&=-\cos i\sin\varepsilon-\sin i\cos\varepsilon\cos\AscendingNode\\
\sin b&=\frac{\cos\varepsilon\sin\AscendingNode}{\sin B}\\
\operatorname{cotang}C&=\frac{\sin i\operatorname{cotang}\varepsilon+\cos i\cos\AscendingNode}{\sin\AscendingNode}\\
\cos c&=\cos i\cos\varepsilon-\sin i\sin\varepsilon\cos\AscendingNode\\
\sin c&=\frac{\sin\varepsilon\sin\AscendingNode}{\sin C}
\end{aligned}
\]

To decide whether $A$, $B$ and $C$ are to be taken in the first two or the
last two quadrants, choose them so that the sines of $a$, $b$ and $c$ are
positive. Thus take $A$ in the first two quadrants if $\cos\AscendingNode$
is positive, and $B$ and $C$ in those same quadrants if $\sin\AscendingNode$
is positive. In the opposite cases take them in the last two quadrants.
Thus $B$ and $C$ lie in the first two or the last two quadrants according
as $\AscendingNode$ lies in the former or the latter.

The fourth, fifth, seventh and eighth of these equations can be put in
a more convenient form by introducing auxiliary angles. This can be done
in two ways.

\textit{First}, put $\frac{\operatorname{tang}i}{\cos\AscendingNode}=\operatorname{tang}E$
and $\operatorname{tang}i\cos\AscendingNode=\operatorname{tang}F$. This gives
\[
\begin{aligned}
\operatorname{cotang}B&=\frac{\sin i\cos(E+\varepsilon)}{\sin\AscendingNode\cos\varepsilon\sin E}=\frac{\cos i\cos(E+\varepsilon)}{\operatorname{tang}\AscendingNode\cos\varepsilon\cos E}\\
\cos b&=-\frac{\cos i\sin(F+\varepsilon)}{\cos F}=-\frac{\sin i\cos\AscendingNode\sin(F+\varepsilon)}{\sin F}\\
\operatorname{cotang}C&=\frac{\sin i\sin(E+\varepsilon)}{\sin\AscendingNode\sin\varepsilon\sin E}=\frac{\cos i\sin(E+\varepsilon)}{\operatorname{tang}\AscendingNode\sin\varepsilon\cos E}\\
\cos c&=\frac{\cos i\cos(F+\varepsilon)}{\cos F}=\frac{\sin i\cos\AscendingNode\cos(F+\varepsilon)}{\sin F}
\end{aligned}
\]

\textit{Second}, put $\frac{\operatorname{tang}\varepsilon}{\cos\AscendingNode}=\operatorname{tang}G$
and $\operatorname{tang}\varepsilon\cos\AscendingNode=\operatorname{tang}H$. This gives
\[
\begin{aligned}
\operatorname{cotang}B&=\frac{\cos(G+i)}{\operatorname{tang}\AscendingNode\cos G}=\frac{\operatorname{tang}\varepsilon\cos(G+i)}{\sin\AscendingNode\sin G}\\
\cos b&=-\frac{\sin\varepsilon\sin(G+i)}{\sin G}=-\frac{\cos\AscendingNode\cos\varepsilon\sin(G+i)}{\cos G}\\
\operatorname{cotang}C&=\frac{\sin(H+i)}{\sin\AscendingNode\operatorname{tang}\varepsilon\cos H}=\frac{\sin(H+i)}{\operatorname{tang}\AscendingNode\sin H}\\
\cos c&=\frac{\cos\varepsilon\cos(H+i)}{\cos H}=\frac{\sin\varepsilon\cos\AscendingNode\cos(H+i)}{\sin H}
\end{aligned}
\]

\SourceLeaf{105}%
It will probably be worth deriving a few more relations between the quantities
$A$, $a$, and so on. The spherical triangle $\mathfrak{x}p\mathfrak{y}$ gives
$\cos\mathfrak{x}\mathfrak{y}=\cos\mathfrak{x}p\cos\mathfrak{y}p+\sin\mathfrak{x}p\sin\mathfrak{y}p\cos\mathfrak{x}p\mathfrak{y}$.
But $\mathfrak{x}\mathfrak{y}=90^0$ and
$\mathfrak{x}p\mathfrak{y}=\mathfrak{x}pP-\mathfrak{y}pP=A-B$. Hence
\[
\cos(A-B)=-\operatorname{cotang}a\operatorname{cotang}b
\]
Likewise, the triangles $\mathfrak{y}p\mathfrak{z}$, $\mathfrak{z}p\mathfrak{x}$ give
\[
\begin{aligned}
\cos(B-C)&=-\operatorname{cotang}b\operatorname{cotang}c\\
\cos(C-A)&=-\operatorname{cotang}c\operatorname{cotang}a
\end{aligned}
\]

Furthermore, in the triangle $\mathfrak{x}p\mathfrak{y}$,
$\cos a=\cos p\mathfrak{y}\mathfrak{x}\sin b$, and in the triangle
$\mathfrak{y}p\mathfrak{z}$, $\sin\mathfrak{z}\mathfrak{y}p=\sin\mathfrak{y}p\mathfrak{z}\sin c$.
Now $\mathfrak{z}\mathfrak{y}p+p\mathfrak{y}\mathfrak{x}=\mathfrak{z}\mathfrak{y}\mathfrak{x}=90^0$,
so we have $\cos a=\sin b\sin c\sin\mathfrak{y}p\mathfrak{z}$ or, since
$\mathfrak{y}p\mathfrak{z}=B-C$,
\[
\sin(B-C)=\frac{\cos a}{\sin b\sin c}
\]
In much the same way we find
\[
\begin{aligned}
\sin(C-A)&=\frac{\cos b}{\sin c\sin a}\\
\sin(A-B)&=\frac{\cos c}{\sin a\sin b}
\end{aligned}
\]

Combining these equations with the preceding ones also gives
\[
\begin{aligned}
\operatorname{cotang}(A-B)&=-\frac{\cos a\cos b}{\cos c}=\operatorname{tang}(F+\varepsilon)\cos a\\
\operatorname{cotang}(B-C)&=-\frac{\cos b\cos c}{\cos a}\\
\operatorname{cotang}(C-A)&=-\frac{\cos c\cos a}{\cos b}=\operatorname{cotang}(F+\varepsilon)\cos a
\end{aligned}
\]
\[
\begin{aligned}
\cos a^2&=\operatorname{cotang}(A-B)\operatorname{cotang}(C-A)\\
\cos b^2&=\operatorname{cotang}(B-C)\operatorname{cotang}(A-B)\\
\cos c^2&=\operatorname{cotang}(C-A)\operatorname{cotang}(B-C)
\end{aligned}
\]
and in a similar way the squares of the sines and tangents of the sides
$a$, $b$, $c$ can be expressed in terms of the angles $A-B$, $B-C$, $C-A$.

\SourceLeaf{106}%
To illustrate the use of these formulas, we shall apply some of them to
Pallas, taking as our basis the latest [VII.] elements of this planet for
1803. We therefore set
\[
\begin{aligned}
i&=34^0\;38'\;1''1\\
\AscendingNode&=172\;28\;13.7\\
\varepsilon&=23\;27\;55.8\quad\text{(mean obliquity according to \textsc{Maskelyne} for 1803).}
\end{aligned}
\]

With these elements the calculation is as follows. The n next to a
logarithm indicates that the corresponding quantity is negative.

\begin{center}\begin{minipage}{\textwidth}\centering
\begin{tabular}{@{}lcl@{}}
$\log\cos i$ & $ .\quad .\quad .$ & $9.9152958$\\
$\log\operatorname{tang}\AscendingNode$ & $ .\quad .$ & $9.1211553\,\mathrm{n}$\\\hline
$\log\operatorname{cotang}A$ & $.$ & $9.0364511$
\end{tabular}\par\smallskip
 thus $A=263^0\;47'\;35''4$
\end{minipage}\end{center}
\begin{center}\begin{minipage}{\textwidth}\centering
\begin{tabular}{@{}lcl@{}}
$\log\cos\AscendingNode$ & $ .\quad .\quad .$ & $9.9962390\,\mathrm{n}$\\
$\log\sin A$ & $ .\quad .\quad .$ & $9.9974467\,\mathrm{n}$\\\hline
$\log\sin a$ & $ .\quad .\quad .$ & $9.9987923$
\end{tabular}
\end{minipage}\end{center}
\begin{center}\begin{minipage}{\textwidth}\centering
\begin{tabular}{@{}lcl@{}}
$\log\sin i$ & $ .\quad .\quad .$ & $9.7545982$\\
$\log\sin\AscendingNode$ & $ .\quad .\quad .$ & $9.1173944$\\\hline
$\log\cos a$ & $ .\quad .\quad .$ & $8.8719926$
\end{tabular}\par\smallskip
 hence $a=85^0\;43'\;44''8$
\end{minipage}\end{center}
\begin{center}\begin{minipage}{\textwidth}\centering
\begin{tabular}{@{}lcl@{}}
$\log\operatorname{tang}i$ & $ .\quad .$ & $9.8393024$\\
$\log\cos\AscendingNode$ & $ .\quad .\quad .$ & $9.9962390\,\mathrm{n}$\quad thus\\\hline
\end{tabular}
\[
\begin{gathered}
\log\operatorname{tang}E\quad .\quad .\quad 9.8430634\,\mathrm{n},\\
E=145^0\;8'\;2''4;\quad E+\varepsilon=168^0\;35'\;58''2\\
\log\operatorname{tang}F\quad .\quad .\quad 9.8355414\,\mathrm{n},\\
F=145\;35\;52.9;\quad F+\varepsilon=169\;3\;48.7
\end{gathered}
\]
\end{minipage}\end{center}
\begin{center}\begin{minipage}{\textwidth}\centering
\begin{tabular}{@{}lcl@{}}
$\log\cos i$ & $ .\quad .\quad .$ & $9.9152958$\\
Compl. $\log\operatorname{tang}\AscendingNode$ & & $0.8788447\,\mathrm{n}$\\
Compl. $\log\cos E$ & & $0.0859260\,\mathrm{n}$\\\hline
$\log\mathrm{const}.$ & $ .\quad .$ & $0.8800665$
\end{tabular}
\end{minipage}\end{center}
\begin{center}\begin{minipage}{\textwidth}\centering
\begin{tabular}{@{}lcl@{}}
$\log\cos(E+\varepsilon)$ & $.$ & $9.9913455\,\mathrm{n}$\\
Compl. $\log\cos\varepsilon$ & & $0.0374886$\\\hline
$\log\operatorname{cotang}B$ & $.$ & $0.9089006\,\mathrm{n}$
\end{tabular}\par\smallskip
 hence $B=172^0\;58'\;7''4$,
\end{minipage}\end{center}
\SourceLeaf{107}%
\begin{center}\begin{minipage}{\textwidth}\centering
\begin{tabular}{@{}lcl@{}}
$\log\mathrm{const}.$ & $ .\quad .$ & $0.8800665$\\
$\log\sin(E+\varepsilon)$ & $.$ & $9.2959318$\\
Compl. $\log\sin\varepsilon$ & & $0.3999023$\\\hline
$\log\operatorname{cotang}C$ & $.$ & $0.5759006$
\end{tabular}\par\smallskip
 $C=14^0\;52'\;12''5$
\end{minipage}\end{center}
\begin{center}\begin{minipage}{\textwidth}\centering
\begin{tabular}{@{}lcl@{}}
$\log\cos\varepsilon$ & $ .\quad .\quad .$ & $9.9625114$\\
$\log\sin\AscendingNode$ & $ .\quad .\quad .$ & $9.1173944$\\
Compl. $\log\sin B$ & & $0.9121791$\\\hline
$\log\sin b$ & $ .\quad .\quad .$ & $9.9920849$
\end{tabular}
\end{minipage}\end{center}
\begin{center}\begin{minipage}{\textwidth}\centering
\begin{tabular}{@{}lcl@{}}
$\log\sin\varepsilon$ & $ .\quad .\quad .$ & $9.6000977$\\
$\log\sin n$ & $ .\quad .\quad .$ & $9.1173944$\\
Compl. $\log\sin C$ & & $0.5906942$\\\hline
$\log\sin c$ & $ .\quad .\quad .$ & $9.3081863$
\end{tabular}
\end{minipage}\end{center}
\begin{center}\begin{minipage}{\textwidth}\centering
\begin{tabular}{@{}lcl@{}}
$\log\cos i$ & $ .\quad .\quad .$ & $9.9152958$\\
$\log\cos F$ & $ .\quad .$ & $9.9165035\,\mathrm{n}$\\\hline
 & & $9.9987923\,\mathrm{n}$
\end{tabular}
\end{minipage}\end{center}
\begin{center}\begin{minipage}{\textwidth}\centering
\begin{tabular}{@{}lcl@{}}
$\log\sin(F+\varepsilon)$ & $.$ & $9.2781142$\\
$\log\cos(F+\varepsilon)$ & $.$ & $9.9920399\,\mathrm{n}$\\\hline
$\log\cos b$ & $ .\quad .\quad .$ & $9.2769065$\\
$\log\cos c$ & $ .\quad .\quad .$ & $9.9908322$
\end{tabular}\par\smallskip
 thus
$\begin{aligned}b&=79^0\;5'\;39''4\\c&=11\;43\;52.8\end{aligned}$
\end{minipage}\end{center}

If only the sines of $a$, $b$, $c$ are wanted, there is no need to compute
their cosines, so the auxiliary angle $F$ can be dispensed with as well.
If we want $a$, $b$, $c$ themselves too, the cosines, which will also be
used later, resolve the ambiguities left by the sines alone. When the sines
are closer to 1, the cosines also give a more precise determination and
at the same time a check on the calculation. For this last purpose it is
also useful that $\frac{\cos i}{\cos F}=\pm\sin a$. The upper sign applies
when $F$ and $A$ are both in the first two quadrants or both in the last
two. The lower sign applies when $F$ is taken in the other half of the
circle from $A$.
\SourceLeaf{108}%
(To see the reason for this, note that in the first case $F$ is equal to
the angle $P\mathfrak{x}p$, and in the second it differs from that angle by
$180^0$.)

The quantities $\varepsilon$, $\AscendingNode$, $i$ are subject to secular
changes. The same will therefore hold for the quantities $A$, $a$, $B$, $b$,
$C$, $c$ that depend on them. If the annual changes in the former are known,
the changes in $A$, $a$, and so on can be calculated from differential formulas
that are easy to derive and that we shall leave aside here. One can also
calculate the values of $A$, $a$, and so on afresh for a more distant epoch,
and deduce their annual changes from these values.

In addition, these quantities undergo periodic changes due to nutation, which
recur with each revolution of the Moon’s nodes. Since the planet’s geocentric
position is wanted relative to the true equator, one should in fact take the
true obliquity of the ecliptic for $\varepsilon$ instead of the mean obliquity,
and for $\AscendingNode$ the distance of the ascending node from the true
equinox instead of the mean equinox. The resulting periodic changes can be
calculated from the same differential formulas as the secular changes and
entered in a table whose argument is the longitude of the Moon’s node. If
many geocentric positions are needed for a period of moderate length and
there is no such table, it is most convenient to calculate the true values
of $A$, $a$, and so on for two epochs at the beginning and end of the period
directly from the true values of $\varepsilon$, $i$, $\AscendingNode$, and
then obtain the values at intervening times by simple interpolation. Over
one year these changes may safely be regarded as uniform.

One could also omit the periodic changes due to nutation altogether and use
the mean values of $A$, $a$, and so on. In that case, for the Earth one would
also have to use the mean obliquity of the ecliptic for $\varepsilon$ and omit
nutation from the longitude $\lambda$ to obtain the distance from the mean
equinox. This gives the planet’s geocentric right ascension and declination
relative to the \textit{mean} equator. Its position relative to the
\textit{true} equator can then be found in the same way as a fixed star’s
mean position is reduced to its apparent position by applying nutation.

\SourceLeaf{109}%
We have only a few remarks on perturbations to add. The perturbations in
latitude, which of course are all that concern us here, are so slight for
all the older planets that they can rightly be disregarded altogether.
They need to be included only for Ceres and Pallas, because of the large
inclination of their orbits to Jupiter’s orbit. There are two ways of doing
this. One can treat the elements that determine the orientation of the
orbit, its inclination and the longitude of its node, as variable and
correct their mean values by periodic equations. Or one can determine
directly how far the planet is forced out of the mean plane of its orbit by
external forces. In the first case one carries these changes over to the
quantities $A$, $a$, and so on, and thus adds other periodic equations
besides those due to nutation. Their arguments will agree with those of the
equations for the inclination and the longitude of the node. This procedure,
however, has not been customary so far. With the second method, by contrast,
the perturbation tables give the perturbation in heliocentric latitude,
which is in fact simply the heliocentric latitude of the planet above the
mean plane of its orbit. Let this be $=\ObservedBeta$, taken as positive toward
the north pole and negative toward the south pole. In the spherical triangle
$\mathfrak{x}pk$, then, the side $pk$ is $=90^0$ in the earlier case and
$=90^0-\ObservedBeta$ here. Hence
\[
x=r\cos\mathfrak{x}k=r(\sin\ObservedBeta\cos a+\cos\ObservedBeta\sin a\sin(t+A))
\]
and likewise
\[
\begin{aligned}
y&=r(\sin\ObservedBeta\cos b+\cos\ObservedBeta\sin b\sin(t+B))\\
z&=r(\sin\ObservedBeta\cos c+\cos\ObservedBeta\sin c\sin(t+C))
\end{aligned}
\]

Since $\ObservedBeta$ here can amount to only a few minutes at most,
one may put $\cos\ObservedBeta=1$ and $\sin\ObservedBeta=\ObservedBeta$.
Thus to allow for the perturbations, one need only take the values of $x$,
$y$, $z$ found without them and add the quantities $\ObservedBeta r\cos a$,
$\ObservedBeta r\cos b$, $\ObservedBeta r\cos c$. Here $\ObservedBeta$ must
be expressed in parts of the radius.

\medskip
[Handwritten remarks on p.\,100.]
\[
\sin a^2=-\frac{\cos(B-C)}{\sin(A-B)\sin(C-A)}
\]
\SourceLeaf{110}%
\[
\operatorname{tang}a^2=-\frac{\cos(B-C)}{\cos(A-B)\cos(C-A)}
\]
\[
\cos a^2+\cos b^2+\cos c^2=1
\]
\[
\begin{aligned}
-\frac{\cos a}{\cos b}&=\operatorname{tang}(B-C)\cos=\frac{\operatorname{cotang}(C-A)}{\cos c}\\
&=\operatorname{tang}\AscendingNode'\ \text{in plano ipsius }Z
\end{aligned}
\]
\[
\begin{aligned}
\sin\AscendingNode'&=\sin b\sin(B-C)=\frac{\cos a}{\sin c}\\
\cos\AscendingNode'&=-\frac{\cos b}{\sin c}
\end{aligned}
\]

\begin{center}Differential changes:\end{center}
\[
2\mathrm{d}a=-\frac{\mathrm{d}A-\mathrm{d}B}{\sin(A-B)\operatorname{cotang}b}-\frac{\mathrm{d}C-\mathrm{d}A}{\sin(C-A)\operatorname{cotang}c}
\]

Conversely, from $a$, $b$, $c$, one can find $\AscendingNode$ and $i$
using the formulas
\[
\begin{aligned}
\cos a&=\sin i\sin\AscendingNode\\
-\cos b\cos\varepsilon-\cos c\sin\varepsilon&=\sin i\cos\AscendingNode\\
-\cos b\sin\varepsilon+\cos c\cos\varepsilon&=\cos i
\end{aligned}
\]
\[
\begin{aligned}
\sin a&=k\\
\cos b&=k\cos\theta\\
\cos c&=k\sin\theta
\end{aligned}
\]
\[
\begin{aligned}
\cos a&=\sin i\sin\AscendingNode\\
-k\cos(\theta-\varepsilon)&=\sin i\cos\AscendingNode\\
k\sin(\theta-\varepsilon)&=\cos i
\end{aligned}
\]
\begin{center}\rule{.25\linewidth}{.4pt}\end{center}
